\section{Rangkuman}
\begin{itemize}
    \item Kurva eliptik $y^2 = x^3 + ax + b$ dengan $4a^3 + 27b^2 \neq 0$; himpunan titiknya bersama titik tak hingga $\mathcal{O}$ membentuk grup abelian $\langle G, + \rangle$.
    \item Penjumlahan titik $P + Q = R$ memakai gradien $m = (y_q - y_p)/(x_q - x_p)$, sedangkan penggandaan $2P$ memakai gradien garis singgung $m = (3x_p^2 + a)/(2y_p)$, dengan $x_r = m^2 - x_p - x_q$ dan $y_r = m(x_p - x_r) - y_p$.
    \item Pada medan berhingga $GF(p)$ semua operasi dilakukan modulo prima $p$, dan pembagian diganti perkalian dengan invers modulo, misalnya $m \equiv (3x_p^2 + a)(2y_p)^{-1} \bmod p$.
    \item Keamanan ECC bersandar pada \textbf{ECDLP}: diberikan $P$ dan $Q = kP$, sangat sukar menemukan skalar $k$ jika $k$ sangat besar.
    \item Kunci privat berupa integer $k \in [1, p-1]$ dan kunci publik berupa titik $Q = kP$; ECDH menghasilkan kunci bersama $K = abB$, sedangkan ECEG mengenkripsi pesan sebagai pasangan titik $[(kB), (P_M + kP_B)]$.
    \item ECC jauh lebih efisien daripada RSA: kunci 256 bit memberi keamanan setara RSA 3072 bit, sehingga cocok untuk \textit{smart card} dan perangkat IoT.
    \item Pesan teks harus dikodekan lebih dulu menjadi titik kurva, misalnya dengan metode Koblitz $x = mk + \text{offset}$ sampai diperoleh $y$ yang valid.
    \item Aritmetika kurva eliptik bertumpu pada aljabar abstrak. Himpunan
          titik pada kurva dengan aturan penjumlahan di atasnya memenuhi kelima
          aksioma grup (Tabel~\ref{tab:aksioma-grup-eliptik}), dan medan
          berhingga menyediakan invers bagi setiap unsur bukan nol
          (Tabel~\ref{tab:contoh-medan}).
    \item Penjumlahan $\Z$ bukan grup karena hanya $1$ dan $-1$ yang memiliki
          invers perkalian, sedangkan $\Z_n$ menjadi medan hanya bila $n$
          prima. Kenyataan itu menentukan modulus mana yang boleh dipakai.
    \item Di dalam $GF(2^m)$ penjumlahan dua polinom tidak lain adalah operasi
          XOR, sedangkan perkalian menuntut reduksi oleh polinom tak
          tereduksi berderajat $m$; sebagai contoh,
          $\code{57} \cdot \code{83} = \code{C1}$ di dalam $GF(2^8)$
          (Tabel~\ref{tab:gf2m-kali}).
    \item Titik di ketakhinggaan $\mathcal{O}$ memainkan dua peran sekaligus:
          elemen netral dan hasil penjumlahan dua titik berlawanan, sehingga
          $P + \mathcal{O} = P$ dan $P + (-P) = \mathcal{O}$
          (Persamaan~\eqref{eq:titik-o}).
    \item Rumus penjumlahan diturunkan dengan menyulihkan persamaan garis ke
          persamaan kurva. Hubungan Vieta memberi $x_p + x_q + x_r = m^2$,
          sehingga $x_r = m^2 - x_p - x_q$
          (Persamaan~\eqref{eq:koordinat-pq}); pada penggandaan akar $x_p$
          muncul dua kali dan $x_r = m^2 - 2x_p$
          (Persamaan~\eqref{eq:koordinat-penggandaan}).
    \item Bila ordinat sebuah titik nol, gradien garis singgungnya tidak
          terdefinisi dan $2P = \mathcal{O}$. Titik semacam itu berorde dua
          dan justru adalah akar-akar polinom kubik kurva.
    \item Perkalian skalar dihitung dengan \textit{double-and-add}, sehingga
          biaya menghitung $kP$ hanya sebanding dengan $\log_2 k$, bukan
          dengan $k$ (Persamaan~\eqref{eq:pelelaran}). Asimetri inilah yang
          membuat ECC praktis sekaligus aman.
    \item Banyaknya titik pada kurva tidak dapat ditentukan hanya dari nilai
          $p$; dalil Hasse hanya membatasinya pada selang
          $p + 1 \pm 2\sqrt{p}$. Karena itu orde grup dan orde titik basis
          harus diperiksa lebih dahulu, dan titik basis harus dipilih yang
          ordonya sama dengan orde grup.
    \item Enumerasi penuh memperlihatkan hal itu secara konkret. Kurva
          $y^2 \equiv x^3 + x + 6 \pmod{11}$ memiliki $12$ titik berhingga
          ditambah $\mathcal{O}$, sehingga orderya $13$
          (Tabel~\ref{tab:enumerasi-gf11} dan Tabel~\ref{tab:pelelaran-gf11}).
    \item Kunci publik ECC adalah sebuah titik, $Q = x \cdot B$
          (Persamaan~\eqref{eq:ecc-kunci-publik}). Dalam ECDH kedua pihak
          menghitung $K = a \cdot P_B = b \cdot P_A = abB$ tanpa pernah
          mengirim kunci privatnya (Persamaan~\eqref{eq:ecdh-kunci}).
    \item ECEG mengenkripsi plainteks $P_M$ menjadi pasangan titik
          $\bigl(kB,\ P_M + kP_B\bigr)$
          (Persamaan~\eqref{eq:eceg-cipher}). Dekripsinya berhasil karena
          $b(kB) = k(bB) = kP_B$, sehingga suku acak pada titik kedua
          terhapus (Persamaan~\eqref{eq:eceg-dekripsi}).
    \item Karena kunci sesaat $k$ dipilih baru untuk setiap pesan, enkripsi
          ECEG bersifat probabilistik: plainteks yang sama menghasilkan
          cipherteks yang berbeda pada setiap pengiriman. Sifat itu tidak
          dimiliki pemetaan karakter langsung, yang selalu memberi titik
          $P_M$ yang sama (Subbagian~\ref{subsec:pemetaan-langsung}).
    \item Metode Koblitz mengubah pesan menjadi titik dengan mencoba
          $x = mk + 1$, lalu $mk + 2$, dan seterusnya sampai ruas kanan kurva
          menjadi kuadrat modulo $p$ (Persamaan~\eqref{eq:koblitz-x});
          pendekodeannya memakai pembulatan ke bawah
          $m = \lfloor (x-1)/k \rfloor$
          (Persamaan~\eqref{eq:koblitz-dekode}).
    \item Algoritma \textit{index calculus} yang menyerang Diffie--Hellman
          dan ElGamal di atas bilangan bulat tidak dapat dipindahkan ke kurva
          eliptik. Serangan terbaik yang diketahui hanyalah rho Pollard, yang
          biayanya sebanding dengan akar orde grup
          (Tabel~\ref{tab:ecdlp-vs-faktorisasi}).
    \item Keunggulan panjang kunci itu hanya berlaku bagi kurva yang dipilih
          dengan benar. Kurva dengan orde yang habis dibagi bilangan kecil
          atau kurva yang dapat ditransfer ke medan berhingga biasa runtuh
          jauh lebih cepat, sehingga kurva baku seperti pada
          Tabel~\ref{tab:kurva-standar} harus dipakai bila memungkinkan.
\end{itemize}
